\loai{Loại 1: Hai vật trao đổi nhiệt – Bài toán thuận}
\begin{ex}
	Một người thợ rèn nhúng một con dao rựa bằng thép có khối lượng $1,1 \ kg$ ở nhiệt độ $850^\circ C$ vào trong bể nước lạnh để làm tăng độ cứng của lưỡi dao. Nước trong bể có thể tích $200$ lít và có nhiệt độ bằng với nhiệt độ ngoài trời là $27^\circ C$. Bỏ qua sự truyền nhiệt cho thành bể và môi trường bên ngoài. Biết nhiệt dung riêng của thép là $460 \ J / kg. K$ của nước là $4180 \ J / kg. K$; khối lượng riêng của nước là $1000\ kg/m^3$. Nhiệt độ của nước khi có sự cân bằng nhiệt là
	\choice
	{$22,4^\circ C$}
	{$55,6^\circ C$}
	{$72,4^\circ C$}
	{\True $27,5^\circ C$}
	\loigiai{
		\begin{itemize}
			\item Gọi nhiệt độ cân bằng là t
			\item Nhiệt lượng do thép truyền cho nước: $Q_{\text{tỏa}}=m_1c_1\Delta t_1=1,1.460.(850-t)$
			\item Nhiệt lượng do nước nhận được từ thép: $Q_{\text{thu}}=m_2c_2\Delta t_2=200.4180.(t-27)$
			\item Khi có sự cân bằng nhiệt: $Q_{\text{tỏa}}=Q_{\text{thu}}\Leftrightarrow t\approx 27,5^\circ C$
		\end{itemize}
	}
	\haidong{6}\end{ex}
%
\begin{ex}%[1P6-1.2G][Loai3]
	Người ta thả miếng đồng có khối lượng $0,5\ kg$ vào bình cách nhiệt chứa 500 g nước. Khi cân bằng nhiệt thì miếng đồng nguội đi từ $80^\circ C$ đến $20^{\circ}C$. Bỏ qua sự trao đổi nhiệt với môi trường và bình chứa. Lấy $c_{Cu}=380\ J/kg.K$, $c_{H_2O}=4190\ J/kg.K$. Nhiệt lượng mà nước nhận được là
	\choice
	{$31600\ J$}
	{\True $11400\ J$}
	{$41100\ J$}
	{$14100\ J$}
	\loigiai{ 
		\begin{itemize}
			\item Nhiệt lượng trao đổi của miếng đồng: $Q_{Cu}=m_{Cu}.c_{Cu}(t-t_1)=-11400$ J.
			\item Theo điều kiện cân bằng nhiệt: $Q_{Cu}+Q_{H_2O}=0\Rightarrow Q_{H_2O}=11400$ J.
			\item Nước nóng lên thêm: 
			\begin{align*}
				Q_{H_2O}=m_{H_2O}.c_{H_2O}.\Delta t\Rightarrow 11400=0,5.4190.\Delta t\Rightarrow \Delta t=5,4^{\circ}C.
			\end{align*}
	\end{itemize}}
	\haidong{6}\end{ex}
\begin{ex}
	Một cốc nhôm khối lượng $100 \ g$ chứa $300 \ g$ nước ở nhiệt độ $20^\circ C$. Người ta thả vào cốc nước một thìa đồng khối lượng $75 \ g$ vừa rút ra khỏi nồi nước sôi ở $100^\circ C$.  Bỏ qua các hao phí nhiệt ra môi trường. Nhiệt dung riêng của nhôm, đồng và nước lần lượt là  $880 \ J / kg. K;\ 380 \ J / kg. K;\ 4190 \ J / kg. K$. Nhiệt độ của nước trong cốc khi có sự cân bằng nhiệt xấp xỉ là
	\choice
	{\True $21,7\doC$}
	{$25,0 \doC$}
	{$28,5 \doC$}
	{$20,2 \doC$}
	\loigiai{
		
		Phương trình cân bằng nhiệt: 
		
		$Q_{\text{cốc}} + Q_{\text{nước}} = Q_{\text{thìa}}$
		
		
		$(m_{\text{cốc}}.c_{\text{cốc}} + m_n.c_n).(t_{cb} – t_1)= m_{\text{thìa}}.c_{\text{thìa}}.(t_2 – t_{cb})$
		
		$\Leftrightarrow [(0,1.880) + (0,3.4,19.103)](t – 20)= 0,075.380.(100 –t)$
		
		Giải ra ta được: $t\approx 21,7^\circ C$.		
	}
	\haidong{7}\end{ex}
\begin{ex}
	Người ta thả một miếng đồng có khối lượng $600\ g$ ở nhiệt độ $100^\circ C$ vào $2,5\ kg$ nước. Nhiệt độ khi có sự cân bằng nhiệt là $30^\circ C$. Lấy nhiệt dung riêng của đồng là $380\ J/kg.K$ và của nước là $4200\ J/kg.K$. Bỏ qua sự trao đổi nhiệt với bình đựng nước và môi trường bên ngoài. Nước nóng lên thêm 
	\choice
	{\True $1,52^\circ C$}
	{$15,2^\circ C$}
	{$1,82^\circ C$}
	{$18,2^\circ C$}
	\loigiai{
		
		Nhiệt lượng tỏa ra của đồng: $Q_{\text{đ}}=m_{\text{đ}} c_{\text{đ}} \Delta t_{\text{đ}}=0,6.380.(100-30)=15960$ J
		
		Nhiệt lượng thu vào của nước $Q_{n}=m_{n} c_{n} \Delta t_{n}=2,5.4200. \Delta t_{n}$
		
		Theo phương trình cân bằng nhiệt: $Q_{\text{đ}}=Q_{n} \Rightarrow 15960=2,5.4200. \Delta t_{n} \Rightarrow \Delta t_{n}=1,52^\circ C$.
		
	}
	\haidong{6}\end{ex}



\begin{ex} \nguon{2}
	Một móng ngựa bằng sắt nặng $1,5 \ kg$, ban đầu ở nhiệt độ $600^\circ C$, rơi vào một cái xô chứa $20 \ kg$ nước ở nhiệt độ $25^\circ C$.  Coi xô cách nhiệt hoàn toàn và giả sử lượng bốc hơi không đáng kể. Biết nhiệt dung riêng của nước và sắt lần lượt là $4200\ J/kg.K$ và $448 \ J / kg.K$. Nhiệt độ của hệ khi cân bằng nhiệt là 
	\choice
	{\True $29,6^\circ C$}
	{$40,5^\circ C$}
	{$15,0^\circ C$}
	{$80,0^\circ C$}
	\loigiai{
		\begin{itemize}
			\item Nhiệt lượng mà móng ngựa ở nhiệt độ cao tỏa ra là
			\begin{align*}
				Q_{\text{tỏa}}=m_{F e} c_{F e}\left(t-t_{F e}\right)
			\end{align*}
			\item Nhiệt lượng mà nước trong xô thu vào là
			\begin{align*}
				Q_{thu}=m_{n} c_{n}\left(t_{n}-t\right)
			\end{align*}
			\item Nhiệt lượng tỏa ra bằng nhiệt lượng thu vào $Q_{\text{tỏa}}=Q_{thu}$
			\item Nhiệt độ cuối cùng của hệ là
			\begin{align*}
				t=\dfrac{m_{n} c_{n} c_{n} t_{n}+m_{F e} c_{F e} t_{F e}}{m_{n} c_{n}+m_{F e} c_{F e}}=\dfrac{20 . 4200 . 25+1,5 . 448 . 600}{1,5 . 448+20 . 4200}=29,6^\circ C
			\end{align*}
		\end{itemize}
	}
\end{ex}

\loai{Loại 2: Hai vật trao đổi nhiệt – Bài toán ngược}

\begin{ex}%[1P6G1-2][Loai4]
	Thả một quả cầu bằng nhôm khối lượng 0,105 kg được đun nóng tới $142^\circ C$ vào một cốc đựng nước ở $20^\circ C$, biết nhiệt độ khi có sự cân bằng nhiệt là $42^\circ C$. Bỏ qua sự truyền nhiệt cho cốc và môi trường bên ngoài. Biết nhiệt dung riêng của nhôm là $880\ J/kg.K$ và của nước là $4200\ J/kg.K$. Khối lượng của nước trong cốc là
	\choice
	{5,01 kg}
	{1,05 kg}
	{\True 0,1 kg}
	{1,5 kg}
	\loigiai{
		\Binhluan{
			\begin{itemize}
				\item Nhiệt lượng nước thu vào để tăng tới $42^\circ C$ là:
				\begin{align*}
					Q_{thu}=m.4200.\left( {42 - 20} \right) = 92400m.
				\end{align*}
				\item Nhiệt lượng quả cầu tỏa ra để giảm tới $42^\circ C$ là: 
				\begin{align*}
					Q_{\text{toả}}=0,105.880.\left( {142 - 42} \right)= 9240\ J.
				\end{align*}
				\item Khi cân bằng nhiệt thì: $Q_{thu}=Q_{\text{toả}}\Rightarrow 9240 = 92400m \to m = 0,1\ kg.$
			\end{itemize}
		}
		{Ta có thể tách ra tính $\sum Q_{thu}=\sum Q_{\text{toả}}$. Khi tính như vậy thì Q luôn dương}
	}
	\haidong{6}\end{ex}


\begin{ex}
	Để xác định nhiệt dung riêng của một kim loại, người ta bỏ vào nhiệt lượng kế chứa $500 \ g$ nước ở nhiệt độ $15^\circ C$ một miếng kim loại có $m=400 \ g$ được đun nóng tới $100^\circ C$. Nhiệt độ khi có sự cân bằng nhiệt là $20^\circ C$. Bỏ qua nhiệt lượng làm nóng nhiệt lượng kế và không khí. Cho nhiệt dung riêng của nước là $4180 \ J /kg.K$. Nhiệt dung riêng của miếng kim loại là
	\choice
	{\True $326,6 \ J / kg. K$}
	{$236,6 \ J / kg. K$}
	{$632,6 \ J / kg. K$}
	{$764,6 \ J / kg. K$}
	\loigiai{
		$c=326,6 \ J / kg. K$
	}
	\haidong{8}\end{ex}

\begin{ex}%[1P6-1.2G][Loai3]
	Để xác định nhiệt độ của một lò nung, người ta đưa vào trong lò một miếng sắt có khối lượng $50\ g$. Khi miếng sắt có nhiệt độ bằng nhiệt độ của lò, người ta lấy ra và thả nó vào một nhiệt lượng kế chứa $900\ g$ nước ở nhiệt độ $17^\circ C$. Khi đó nhiệt độ của nước tăng lên đến $23^\circ C$.  Biết nhiệt dung riêng của sắt là $478\ J/kg.K$, của nước là $4180\ J/kg.K$. Bỏ qua sự trao đổi nhiệt với môi trường và với nhiệt lượng kế. Nhiệt độ của lò xấp xỉ bằng
	\choice
	{$813^\circ C$}
	{$796^\circ C$}
	{$990^\circ C$}
	{\True $967^\circ C$}
	\loigiai{
		\begin{itemize}
			\item Nhiệt lượng do nước thu vào là:
			\begin{align*}
				Q_{thu}=0,9.4180.\left( {23 - 17} \right) = 22572\ J.
			\end{align*}
			\item Nhiệt lượng do sắt tỏa ra là:
			\begin{align*}
				Q_{\text{tỏa}}=0,05.478.\left( {t - 23} \right) = 23,9t - 549,7.
			\end{align*}
			\item Khi cân bằng nhiệt thì:
			\begin{align*}
				Q_{thu}=Q_{\text{tỏa}}\to 22572 = 23,9t - 549,7\to t = 967^\circ C.
			\end{align*}
		\end{itemize}
	}
	\haidong{6}\end{ex}
%




\begin{ex} \nguon{3}
	Người ta thả một miếng đồng có khối lượng $m_1=0,2\ kg$ đã được đốt nóng đến nhiệt độ $t_1$ vào một nhiệt lượng kế chứa $m_2=0,28\ kg$ nước ở nhiệt độ $t_2=20^\circ C$. Nhiệt độ khi có cân bằng nhiệt là $t_3=80^\circ C$. Biết nhiệt dung riêng của đồng và nước lần lượt là $c_1=400\ J/kg . K$, $c_2=4200\ J/kg . K$. Bỏ qua sự trao đổi nhiệt với nhiệt lượng kế và với môi trường. Nhiệt độ ban đầu $t_1$ của đồng bằng bao nhiêu độ C (làm tròn kết quả  đến chữ số hàng đơn vị)?
	\shortans[oly]{962 } 
	\loigiai{
		Độ lớn nhiệt lượng đồng tỏa ra bằng độ lớn nhiệt lượng mà nước thu vào:
		$$\left|Q_1\right|=\left|Q_2\right| \Leftrightarrow m_1 c_1\left(t_1-t_3\right)=m_2 c_2\left(t_3-t_2\right) \Leftrightarrow 0,2 . 400 . \left(t_1-80\right)=0,28 . 4200 . (80-20) \Rightarrow t_1\approx 962^\circ C$$
	}
\end{ex}



\begin{ex} \nguon{4}
	Một quả cầu nhôm có khối lượng $0,15 \ kg$ được đun nóng tới nhiệt độ $t_1=$ $150^\circ C$ rồi thả vào chậu nước ở nhiệt độ $t_2=25^\circ C$. Sau một thời gian nhiệt độ của cả hệ là $t=30^\circ C$. Xem nhiệt lượng chỉ trao đổi giữa nước và quả cầu. Biết nhiệt dung riêng của nước là $4180 \ J/kg.K$ và của nhôm là $880 \ J/kg.K$. Khối lượng của nước là
	\choice
	{0,578 kg}
	{0,857 kg}
	{\True 0,758 kg}
	{0,785 kg}
	\loigiai{
		\begin{itemize}
			\item Ta có: $T_1=t_1+273=423 \ K, \ T_2=t_2+273=298 \ K, \ T=t+273=303 \ K$
			\item Nhiệt lượng quả cầu nhôm tỏa ra khi hạ nhiệt từ $150^\circ C$ xuống $30^\circ C$:
			\item $Q_1=m_1 c_1\left(\ T_1-T\right)=0,15.880.(423-303)=15840 \ J$
			\item Nhiệt lượng nước thu vào để nó tăng nhiệt độ từ $25^\circ C$ đến $30^\circ C$:
			\item $Q_2=m_2. c_2\left(\ T-T_2\right)=m_2 . 4180.(303-298)=20900 \ m_2 \ J$
			\item Theo phương trình cân bằng nhiệt: $Q_1=Q_2$
			\item Tính được: $m_2=0,758 \ kg$.
		\end{itemize}
	}
\end{ex}




%	 

\begin{ex} \nguon{4}
	Một người muốn làm tăng nhiệt độ của $20 \ kg$ nước ở $20^\circ C$ tăng lên đến $40^\circ C$. Bỏ qua sự trao đổi nhiệt của nước với môi trường và bình chứa. Khối lượng của nước ở $100^\circ C$ cần dùng để pha vào lượng nước trên là  
	\choice
	{$9,67\ kg$}
	{$6,76\ kg$}
	{$7,67\ kg$}
	{\True $6,67\ kg$}	
	\loigiai{
		\begin{itemize}
			\item $T_1=t_1+273=293 \ K, \ T_2=t_2+273=313 \ K, \ T_3=t_3+273=373 \ K$
			\item Nhiệt lượng $20 \ kg$ nước thu vào để tăng nhiệt độ từ $20^\circ C$ đến $40^\circ C$
			
			$Q_1=m_1.c. \left(T_2-T_1\right)=20. c.(313-293)=400 c$
			\item Nhiệt lượng do khối nước nóng tỏa ra khi hạ nhiệt từ $100^\circ C$ xuống $40^\circ C$.
			
			$Q_2=m_2.c. \left(T_3-T_2\right)=m_2$.c. $(373-313)=60 \ m_2 c$
			\item Từ $Q_1=Q_2$, tính được $m_2=6,67 \ kg$.
		\end{itemize}
	}
\end{ex}

\loai{Loại 3: Ba vật trao đổi nhiệt}
\begin{ex}%[1P6-1.2T][Loai6] 
	Một bình nhôm khối lượng $0,5\ kg$ chứa $0,118\ kg$ nước ở nhiệt độ $20^\circ C$. Người ta thả vào bình một miếng sắt khối lượng $0,2\ kg$ đã được nung nóng tới $75^\circ C$. Bỏ qua sự truyền nhiệt ra môi trường bên ngoài. Nhiệt dung riêng của nhôm là $896\ J/kg.K$; của nước là $4180\ J/kg.K$, của sắt là $460\ J/kg.K$. Nhiệt độ của nước khi bắt đầu có sự cân bằng nhiệt là
	\choice
	{$23,6^\circ C$}
	{$29,4^\circ C$}
	{$25,2^\circ C$}
	{\True $24,9^\circ C$}
	\loigiai{
		\begin{itemize} 
			\item Nhiệt lượng trao đổi của bình nhôm: $Q_{Al} = m_{Al}c_{Al}(t - t_1)$.
			\item Nhiệt lượng trao đổi của nước: $Q_{H_2O} = m_{H_2O}c_{H_2O}(t - t_1)$.
			\item Nhiệt lượng trao đổi của sắt: $Q_{Fe} = m_{Fe}c_{Fe}(t - t_2)$.
			\item Khi có cân bằng nhiệt: 
			\begin{align*}
				Q_{Al} +Q_{H_2O}+Q_{Fe}=0 \Rightarrow t \approx 24,9^{\circ}C.
			\end{align*}
	\end{itemize}}
	\haidong{7}\end{ex}
	\begin{ex}%[1P6-1.2T][Loai8] 
	Một nhiệt lượng kế bằng đồng khối lượng có $100\ g$ chứa $375\ g$ nước ở nhiệt độ $25^\circ C$. Cho vào nhiệt lượng kế một vật bằng kim loại khối lượng $400\ g$ ở $90^\circ C$. Biết nhiệt độ khi có sự cân bằng nhiệt là $30^\circ C$. Bỏ qua sự trao đổi nhiệt với môi trường. Cho biết nhiệt dung riêng của đồng là $380\ J/kg.K$ và của nước là $4200\ J/kg.K$. Nhiệt dung riêng của miếng kim loại là
	\choice
	{$633\ J/kg. K$}
	{\True $336\ J/kg. K$}
	{$362\ J/kg. K$}
	{$880\ J/kg. K$}
	\loigiai{
		\Binhluan{
			\begin{itemize} 
				\item Nhiệt lượng trao đổi của nhiệt lượng kế: $Q_{NLK} = m_{Cu}c_{Cu}(t - t_1)$.
				\item Nhiệt lượng trao đổi của nước: $Q_{H_2O} = m_{H_2O}c_{H_2O}(t - t_1)$.
				\item Nhiệt lượng trao đổi của kim loại: $Q_{Kl} = m_{Kl}c_{Kl}(t - t_2)$.
				\item Khi có cân bằng nhiệt: 
				$Q_{NLK} +Q_{H_2O}+Q_{Kl}=0 \Rightarrow c_{Kl} = 336\ J/kg.K.$
			\end{itemize}
		}
		{Đặt ẩn số cần tìm là x và dùng SHIFT SOLVE}
	}
	\haidong{8}\end{ex}





\begin{ex}%[1P6-1.2T][Loai6]
	Trộn ba chất lỏng không tác dụng hoá học lẫn nhau. Biết $m_1=1\ kg$, $m_2=10\ kg$, $m_3=5\ kg$, $t_1=6^{\circ}C$, $t_2=-40^{\circ}C$, $t_3=60^{\circ}C$, $c_1=2\ kJ/kg.K$, $c_2=4\ kJ/kg.K$, $c_3=2\ kJ/kg.K$. Coi nhiệt chỉ trao đổi giữa các chất lỏng. Nhiệt độ khi cân bằng là
	\choice
	{$t=-11^{\circ}C$}
	{\True $t=-19^{\circ}C$}
	{$t=-29^{\circ}C$}
	{$t=19^{\circ}C$}
	\loigiai{ 
		\begin{itemize}
			\item Theo điều kiện cân bằng nhiệt: 
			\begin{align*}
				Q_1+Q_2+Q_3=0\Rightarrow =m_1.c_1.(t-t_1)+m_2.c_2.(t-t_2)+m_3.c_3.(t-t_3)= 0 \Rightarrow t=-19^{\circ}C.
			\end{align*}
		\end{itemize}
	}
	\haidong{10}\end{ex}




\begin{ex} \nguon{2}
	Một nhiệt lượng kế bằng đồng có khối lượng $0,1 \ kg$ chứa $0,5 \ kg$ nước ở $20^\circ C$. Người ta thả vào nhiệt lượng kế nói trên một thỏi đồng có khối lượng $0,2 \ kg$ đã được đun nóng đến $200^\circ C$. Cho biết nhiệt dung riêng của đồng là $380 \ J /kg.K$ và của nước là $4200 \ J /kg.K$. Bỏ qua sự trao đổi nhiệt với môi trường. Nhiệt độ cuối cùng của hệ là
	\choice
	{$28,2^\circ C$}
	{$28^\circ C$}
	{$27,4^\circ C$}
	{\True $26,2^\circ C$}
	\loigiai{
		\begin{itemize}
			\item Nhiệt lượng mà nhiệt lượng kế thu vào
			\begin{align*}
				Q_1=m_1 c_1\left(t_2-t_1\right)=0,1 . 380.\left(t_2-20\right)=38\left(t_2-20\right)
			\end{align*}
			\item Nhiệt lượng nước thu vào
			\begin{align*}
				Q_2=m_2 . c_2\left(t_2-t_1\right)=0,5 . 4200 .\left(t_2-20\right)=2100.\left(t_2-20\right)
			\end{align*}
			\item Nhiệt lượng đồng tỏa ra
			\begin{align*}
				Q_3=m_3 . c_1 .\left(t_3-t_2\right)=1,2 . 380 .\left(200-t_2\right)=76 .\left(200-t_2\right)
			\end{align*}
			\item Theo phương trình cân bằng nhiệt ta có: $Q=Q_1+Q_2$
			\begin{align*}
				& \Rightarrow 38. t_2-760+2100 t_2-4200=15200-t_2 
				\Rightarrow t_2\approx 26,2^\circ C
			\end{align*}
		\end{itemize}
	}
\end{ex}






\begin{ex} \nguon{4}
	Người ta trộn một hỗn hợp từ ba chất lỏng không có tác dụng hóa học với nhau có khối lượng $m_1=1 \ kg, \ m_2=2 \ kg, \ m_3=3 \ kg$, các chất lỏng có nhiệt độ lần lượt là $t_1=10^\circ C,\ t_2=20^\circ C$ và $t_3=30^\circ C$. Biết nhiệt dung riêng của các chất lỏng lần lượt là $c_1=2500 \ J/kg.K$, $c_2=4200 \ J/kg.K$ và $c_3=3900 \ J/kg.K$.  Bỏ qua sự trao đổi nhiệt với môi trường bên ngoài. Nhiệt độ của hỗn hợp khi nó đạt trạng thái cân bằng nhiệt là
	\choice
	{\True $24,1^\circ C$}
	{$14,1^\circ C$}
	{$34,1^\circ C$}
	{$44,1^\circ C$}
	\loigiai{
		Khi hỗn hợp đạt trạng thái cân bằng nhiệt:
		\begin{align*}
		&m_1 c_1\left(t_1-t\right)+m_2 c_2\left(t_1-t\right)+m_3 c_3\left(t_3-t\right)=0\\
		&t=\dfrac{m_1 c_1 t_1+m_2 c_2 t_2+m_3 c_3 t_3}{\ m_1 c_1+m_2 c_2+m_3 c_3}=24,1^\circ C
		\end{align*}
	}
\end{ex}


\begin{ex} \nguon{4}
	Người ta cho vòi nước nóng $t_{n}=80^\circ C$ và vòi nước lạnh $t_{\ell}=10^\circ C$ đồng thời chảy vào bể đã có sẵn 100 lít nước ở nhiệt độ $t_0=30^\circ C$. Bỏ qua sự trao đổi nhiệt với môi trường bên ngoài.  Cho biết lưu lượng của mỗi vòi là 20 lít/phút. Phải mở hai vòi trong bao nhiêu phút thì thu được nước có nhiệt độ $t=40^\circ C$?
	\shortans[oly]{5 }
	\loigiai{
		\begin{itemize}
			\item $T_{n}=353 \ K, \ T_{\ell}=283 \ K, \ T_0=303 \ K, \ T=313 \ K, \ m_0=100 \ kg$.
			\item Gọi $m$ là lượng nước mỗi vòi chảy vào bể.
			\item Khi hệ cân bằng, nước nóng sẽ tỏa nhiệt, nước trong bể và nước lạnh sẽ thu nhiệt lượng: $mc\left(T_{n}-T\right)=m_0 c\left(T-T_0\right)+mc\left(T-T_{l}\right)$
			\item Tính được: $m=100 \ kg$, thể tích $V=100$ lít.
			\item Thời gian cần thiết: $\Delta t=\dfrac{100}{20}=5$ phút.
		\end{itemize}
	}
\end{ex}

\loai{Loại 4: Bốn vật trao đổi nhiệt}

\begin{ex} \nguon{2}
	Một miếng ghép có tổng khối lượng $m=50 \ g$ từ chì và kẽm ở $136^\circ C$ được thả vào một nhiệt lượng kế. Nhiệt lượng kế có  nhiệt dung $30\ J/K$ và chứa $100 \ g$ nước ở $14^\circ C$. Nhiệt độ khi cân bằng nhiệt là $18^\circ C$.  Biết nhiệt dung riêng của nước là $c_0=4200 \ J /kg.K$, của chì là $c_1=130 \ J /kg.K$, của kẽm là $c_2=380 \ J /kg.K$. Bỏ qua sự trao đổi nhiệt với môi trường bên ngoài. Khối lượng chì là
	\choice
	{\True $15 \ g$}
	{$25 \ g$}
	{$5 \ g$}
	{$35 \ g$}
	\loigiai{
		\begin{itemize}
			\item Phương trình cân bằng nhiệt cho hệ: 
			\begin{align*}
				& Q_1+Q_2+Q_3+Q_4=0\\
				\Leftrightarrow \ & c_1 m_1\left(t-t_1\right)+c_2 m_2\left(t-t_2\right)+c_3 m_3\left(t-t_3\right)+c_0 m_0\left(t-t_4\right)=0
			\end{align*}
			\item Trong đó: 
			\begin{itemize}
				\item $t_1=t_2=136^\circ C;\ t_3=t_{4}=14^\circ C;\ c_3 m_3=30 \ J / K;\ t=18^\circ C$;
				\item $c_1=130 \ J / g . K;\ c_2=380 \ J / kg . K;\ c_0=4200 \ J / kg . K;\ m_0=100 \ g.$ 
			\end{itemize}
			\begin{align*}	
				& \Rightarrow 130 . m_1.(18-136)+380 . (0,05-m_1).(18-136)+30.(18-14)+4200 . 0,1.(18-14)=0  \qquad \text {(1)}
			\end{align*}
			
			\item Giải  (1) ta được: $m_1\approx 15 \ g$.
		\end{itemize}
	}
\end{ex}

%



\loai{Loại 5. Nâng cao}
\begin{ex} \nguon{4}
	Người ta thả một chai sữa của trẻ em có nhiệt độ $t_0$ vào phích đựng nước ở nhiệt độ $40^\circ C$. Khi đạt cân bằng nhiệt, chai sữa nóng tới nhiệt độ $36^\circ C$, người ta lấy chai sữa này ra. Ngay sau đó, tiếp tục thả vào phích một chai sữa khác giống như chai sữa trên có cùng nhiệt độ $t_0$, khi cân bằng chai sữa này nóng tới nhiệt độ $33^\circ C$. Bỏ qua sự trao đổi nhiệt với môi trường bên ngoài. Giá trị của $t_0$ là bao nhiêu độ C?
	\shortans[oly]{24 }
	\loigiai{
		\begin{itemize}
			\item Ta có: 
			\begin{align*}
			&m_{s} c_{s}\left(36-t_0\right)=4m_{n} c_{n}\\
			&m_{s} c_{s}\left(33-t_0\right)=3m_{n} c_{n}
			\end{align*}
			\item Tính được $t_0=24^\circ C$
		\end{itemize}
	}
\end{ex}

\begin{ex} \nguon{2}
	Một bình hình trụ lớn chứa một lượng chất lỏng ở nhiệt độ $60^\circ C$. Nếu thả vào bình một quả cầu kim loại có khối lượng M ở nhiệt độ $100^\circ C$ thì nhiệt độ cân bằng của hệ là $80^\circ C$. Nếu thả quả cầu cùng chất như trên nhưng có khối lượng $2M$ ở nhiệt độ $t^\circ C$ thì nhiệt độ cân bằng của hệ là $50^\circ C$. Coi như quá trình trao đổi nhiệt chỉ xảy ra giữa chất lỏng và quả cầu; quả cầu ngập hoàn toàn trong chất lỏng. Biết rằng nước không bị tràn ra ngoài. Giá trị $t$ bằng bao nhiêu độ C?
	\shortans[oly]{45} 
	\loigiai{
		\begin{itemize}
			\item Gọi $m_0$ là khối lượng chất lỏng trong bình.\\
			$c_0, c$ lần lượt là nhiệt dung riêng của chất lỏng, chất làm quả cầu.
			\item Khi thả quả cầu khối lượng $M$, ta có phương trình cân bằng nhiệt:
			\begin{align*}
				m_0 .c_0.(80-60)=M. c.(100-80) \Rightarrow m_0. c_0=M. c
			\end{align*}
			\item Khi thả quả cầu khối lượng $2 M$, ta có phương trình cân bằng nhiệt:
			\begin{align*}
				m_0. c_0.(60-50)=2. M .c.(50-t) \Rightarrow 60-50=2.(50-t) \Rightarrow t=45^\circ C
			\end{align*}
		\end{itemize}
	}
\end{ex}

\begin{ex} \nguon{2}
	Dùng một ca múc nước ở thùng chứa nước $A$ có nhiệt độ $80^{\circ} C$ và ở thùng chứa nước $B$ có nhiệt độ $20^{\circ} C$ rồi đổ vào thùng chứa nước $C$. Biết rằng trước khi đổ, trong thùng chứa nước $C$ đã có sẵn một lượng nước ở nhiệt độ $40^{\circ} C$ và và bằng tổng số ca nước vừa đổ thêm và nó. Bỏ qua trao đổi nhiệt với môi trường, với bình chứa và ca múc. Thể tích nước lấy ở mỗi lần múc là như nhau. Số ca nước lấy từ thùng A gấp mấy lần số ca nước lấy từ thùng B để nhiệt độ nước ở thùng C là $50^\circ C$?
	\shortans[oly]{ 2} 
%	\choice
%	{$3$}
%	{\True $2$}
%	{$2,5$}
%	{$4$}
	\loigiai{
		\begin{itemize}
			\item Gọi m là khối lượng nước múc từ ca
			$\mathrm{n}_{\mathrm{a}}$. 
			\item $\mathrm{n}_b$ lần lượt là ca múc từ thùng $\mathrm{A}$ và thùng $\mathrm{B}$
			\item Ta có phương trình cân bằng nhiệt:
			$$
			\begin{aligned}
				& \mathrm{Q}_{\mathrm{A}}+\mathrm{Q}_{\mathrm{B}}+\mathrm{Q}_{C}=0 \\
				& \Leftrightarrow m_A c\left(t_A-t\right)+m_B c\left(t_B-t\right)+m_C c\left(t_C-t\right)=0 \\
				& \Rightarrow n_A=2 n_B
			\end{aligned}
			$$
		\end{itemize}
	}
\end{ex}

\begin{ex} \nguon{2}
	Một bình cách nhiệt nhẹ chứa nước ở nhiệt độ $20^\circ C$. Người ta lần lượt thả vào bình này những quả cầu giống nhau đã được đốt nóng đến $100^\circ C$. Sau khi thả quả cầu thứ nhất thì nhiệt độ của nước trong bình khi cân bằng nhiệt là $40^\circ C$. Biết nhiệt dung riêng của nước là $4200 \ J /kg.K$. Bỏ qua sự trao đổi nhiệt với bình và môi trường. Giả thiết nước không bị tràn ra ngoài và không tính đến sự bay hơi của nước. 
	Để nhiệt độ trong bình khi cân bằng nhiệt là $90^\circ C$ thì phải thả vào bình bao nhiêu quả cầu?
	\shortans[oly]{ 21}
%	\choice
%	{\True $21$ quả cầu}
%	{$25$ quả cầu}
%	{$15$ quả cầu}
%	{$10$ quả cầu}
	\loigiai{
		\begin{itemize}
			\item Cân bằng nhiệt lần đầu: $$m_1.4200.(40-20)=m_2c_2.(100-40)$$
			\item Gọi $n$ là số quả cầu cần thả vào.
			\item Cân bằng nhiệt lần sau:			
				$4m_1.4200.(90-20)=n.m_2.c.10$4
		\item 	Vậy phải thả tổng cộng 21 quả cầu.
		\end{itemize}
	}
\end{ex}
\begin{ex}
	Ba bình đựng ba chất lỏng khác nhau và không gây tác dụng hóa học với nhau. Nhiệt độ của ba bình lần lượt là $t_1=30^{\circ}C,t_2=10^{\circ}C$ và $t_3=45^{\circ}C$. Nếu đổ một nửa chất lỏng ở bình 1 sang bình 2 thì nhiệt độ của hỗn hợp khi cân bằng nhiệt độ là $t_{12}=16^{\circ}C$. Còn nếu đổ một nửa chất lỏng ở bình 1 sang bình 3 thì nhiệt độ của hỗn hợp khi cân bằng nhiệt là $t_{13}=35^{\circ}C$. Xem như chỉ có các chất lỏng trao đổi nhiệt với nhau.
	\choiceTFt
	{Nếu đổ một nửa chất lỏng ở bình 1 sang bình 2 thì chất lỏng bình 1 nhận nhiệt lượng từ bình 2}
	{\True Nếu đổ một nửa chất lỏng ở bình 1 sang bình 3 thì chất lỏng ở bình 1 nhận nhiệt lượng từ bình 3}
	{\True Khi đổ cả chất lỏng ở bình 1 vào chất lỏng ở bình 2 thì nhiệt độ hỗn hợp chất lỏng khi có sự cân bằng nhiệt là $19,2^{\circ}C$}
	{\True Nếu đổ cả ba chất lỏng vào một bình thì nhiệt độ của hỗn hợp chất lỏng khi có sự cân bằng nhiệt là $21,9^{\circ}C$}
	\loigiai{\begin{itemize}
			\item Sai. Nếu đổ một nửa chất lỏng ở bình 1 sang bình 2 thì chất lỏng bình 2 nhận nhiệt lượng từ bình 1.
			\item Đúng.
			\item Đúng. Khi đổ một nửa chất lỏng ở bình 1 vào chất lỏng ở bình 2, ta có phương trình cân bằng nhiệt: 
			\begin{align*}
				Q_1+Q_2=0\Rightarrow\dfrac{m}{2} c_1(16-30)+m c_2(16-10)=0\Rightarrow 14 c_1=12 c_2\Rightarrow c_1=\dfrac{6}{7} c_2
			\end{align*}
			Khi đổ cả chất lỏng ở bình 1 vào chất lỏng ở bình 2, ta có phương trình cân bằng nhiệt: 
			\begin{align*}
				Q_1{}^{\prime}+Q_2{}^{\prime}=0\Rightarrow m c_1\left(t_{12}{}^{\prime}-30\right)+m c_2\left(t_{12}{}^{\prime}-10\right)=0
				\Rightarrow m\dfrac{6}{7} c_2\left(t_{12}^{\prime}-30\right)+m c_2\left(t_{12}^{\prime}-10\right)=0\Rightarrow t_{12}^{\prime}=19,2^{\circ}C
			\end{align*}
			\item Đúng. Khi đổ một nửa chất lỏng ở bình 1 vào chất lỏng ở bình 3, ta có phương trình cân bằng nhiệt: 
			\begin{align*}
				Q_1+Q_3=0
				\Rightarrow\dfrac{m}{2} c_1(35-30)+m c_3(35-45)=0\Rightarrow 5 c_1=20 c_3\Rightarrow c_1=4 c_3 
			\end{align*}
			Khi đổ chất lỏng ở cả 3 bình vào 1, ta có phương trình cân bằng nhiệt:
			\begin{align*}
				Q_1^{\prime}+Q_2^{\prime}+Q_3^{\prime}=0
				&\Rightarrow m c_1\left(t^{\prime}-30\right)+m c_2\left(t^{\prime}-10\right)+m c_3\left(t^{\prime}-45\right)=0\\
				&\Rightarrow m c_1\left(t^{\prime}-30\right)+m\dfrac{7}{6} c_1\left(t^{\prime}-10\right)+m\dfrac{1}{4} c_1\left(t^{\prime}-45\right)=0\Rightarrow t^{\prime}=21,9^{\circ} C
			\end{align*}
		\end{itemize}
	}
\end{ex}
\begin{ex} \nguon{2}
	Trên bàn có ba bình giống nhau đựng các lượng nước như nhau ở cùng nhiệt độ. Đổ $m$ gam nước nóng vào bình thứ nhất, khi có cân bằng nhiệt thì múc $m$ gam nước từ bình thứ nhất đổ vào bình thứ hai. Sau đó múc $m$ gam nước từ bình thứ hai đã cân bằng nhiệt đổ vào bình thứ ba. Độ tăng nhiệt độ của nước ở bình thứ nhất và thứ hai lần lượt là $\Delta t_1$ $=20^\circ C$ và $\Delta t_2=16^\circ C$. Bỏ qua sự trao đổi nhiệt với bình và môi trường. Độ tăng nhiệt độ $\Delta t_3$ của nước ở bình thứ ba bằng bao nhiêu độ C (làm tròn kết quả đến chữ số hàng phần mười)?
	\shortans[oly]{12,8 }

	\loigiai{
		\begin{itemize}
			\item Gọi nhiệt độ ban đầu của các bình chứa nước là $t^\circ C$\\
			Khối lượng nước trong mỗi bình là $m_1\ (kg)$
			\item Nhiệt độ sau khi cân bằng lần thứ nhất là:
			\begin{align*}
				t+\Delta t_1=t+20\left({ }^\circ C\right)
			\end{align*}
			\item Nhiệt độ sau khi cân bằng lần thứ hai là:
			\begin{align*}
				t+\Delta t_2=t+16\left({ }^\circ C\right)
			\end{align*}
			\item Nhiệt độ sau khi cân bằng lần thứ ba là:
			\begin{align*}
				t+\Delta t_3\left({ }^\circ C\right)
			\end{align*}
			\item Áp dụng phương trình cân bằng nhiệt cho lần thứ hai, ta có:
			\begin{align*}
				& Q_{\text {thu}\ 2}=Q_{\text {tỏa}\ 2} \\
				\Leftrightarrow \ & m_1 . c . \Delta t_2=m . c .[(t+20)-(t+16)] \\
				\Leftrightarrow \ & m_1 . 16=m . 4 \\
				\Leftrightarrow \ & m=4 . m_1
			\end{align*}
			\item Áp dụng phương trình cân bằng nhiệt cho lần thứ ba, ta có:
			\begin{align*}
				& Q_{\text {thu}\ 3}=Q_{\text {tỏa}\ 3} \\
				\Leftrightarrow \ & m_1 . c . \Delta t_3=m . c .\left[(t+16)-\left(t+\Delta t_3\right)\right] \\
				\Leftrightarrow \ & m_1 . \Delta t_3=4 m_1 .\left(16-\Delta t_3\right) \\
				\Leftrightarrow \ & \Delta t_3=64-4 \Delta t_3 \\
				\Leftrightarrow \ & 5 . \Delta t_3=64 \\
				\Leftrightarrow \ & \Delta t_3=12,8^\circ C 
			\end{align*}
		\end{itemize}
	}
\end{ex}



\begin{ex} \nguon{2}
	\impicinpar{Một bình hình trụ chứa $800 \ g$ nước ở $20^\circ C$. Bình được đặt trên một mặt bàn nằm ngang (hình bên). Người ta thả nhẹ vào bình một thanh sắt đặc đồng chất, hình trụ có thể tích $V_1=300\ cm^3$ ở $100^\circ C$. Biết nhiệt dung riêng của sắt và nước lần lượt là $c_1=460\ J / kg . K,\ c_2=4200 \ J / kg . K$. khối lượng riêng của sắt là $\rho_1=7800 \ kg / m^3$. 	 Bỏ qua sự trao đổi nhiệt giữa nước và môi trường. Bỏ qua sự hóa hơi của nước và sự nở vì nhiệt của các vật. Cho biết hiệu suất của quá trình là $H=80 \%$ (với $H=\dfrac{Q_2}{Q_1}.100 \%$, trong đó $Q_2$ là nhiệt lượng thu vào của nước và $Q_1$ là nhiệt lượng tỏa ra của thanh sắt) và cho rằng khi có sự cân bằng nhiệt, nhiệt độ tại mọi điểm của nước và của sắt đều bằng nhau. Nhiệt độ của hệ nước và thanh sắt khi có sự cân bằng nhiệt là bao nhiêu độ C (làm tròn kết quả đến chữ số hàng phần mười)?}
	{\includegraphics[width=6.5cm]{images/binhnuocthanhsat}}
	\shortans[oly]{ 36,3} 
	%	\choice
	%	{\True $36,32^\circ C$}
	%	{$35,15^\circ C$}
	%	{$34,34^\circ C$}
	%	{$34,32^\circ C$}
	\loigiai{
		Gọi t là nhiệt độ cân bằng\\
		Nhiệt lượng $800 \ g$ nước thu:
		\begin{align*}
			Q_2=m_2 c_2 \Delta t_2=0,8.4200(t-20)=3360 t-67200
		\end{align*}
		Nhiệt lượng thanh sắt tỏa:
		\begin{align*}
			& Q_1=m_1 c_1 \Delta t_1=\rho_1 V_1 c_1 \Delta t_1 \\
			\Rightarrow \ & Q_1=7800.300 .10^{-6} .460.(100-t)=107640-1076,4 t 
		\end{align*}
		Hiệu suất của quá trình là $80 \%=0,8$ nên:
		\begin{align*}
			\dfrac{3360 t-67200}{107640-1076,4 t}=0,8 \Rightarrow t=36,32^\circ C
		\end{align*}
	}
\end{ex}

















\begin{ex} \nguon{4}
	Người ta nung nóng miếng đồng có khối lượng $100 \ g$ đến nhiệt độ $650^\circ C$ rồi thả vào cốc nước có thể tích 1 lít đang có nhiệt độ $30^\circ C$. Giả sử cốc nước được làm từ vật liệu cách nhiệt rất tốt và bỏ qua sự trao đổi nhiệt với môi trường bên ngoài. Biết khối lượng riêng của nước là $1000 \ kg / m^3$ và nhiệt dung riêng của đồng và của nước lần lượt là $c_1=380 \ J/kg.K$ và $c_2=4180 \ J/kg.K$.
	\choiceTF[t]
	{\True Đã có quá trình truyền nhiệt từ miếng đồng sang nước}
	{\True Khi hệ đã cân bằng nhiệt, so với ban đầu thì nội năng của miếng đồng đã giảm xuống, còn của nước tăng lên}
	{Khi hệ đã cân bằng nhiệt, nhiệt độ của của nước trong cốc là $40,5^\circ C$}
	{Độ biến thiên nội năng của miếng đồng xấp xỉ là $23348 \ J$}
	\loigiai{
		\begin{itemchoice}
			\itemch 
			\itemch 
			\itemch Sai. Ta có: $m_1 c_1\left(t_1-t\right)=m_2 c_2\left(t-t_2\right) \Rightarrow$ Nhiệt độ lúc sau của hệ là $35,6^\circ C$.
			\itemch Đúng. Áp dụng công thức. $\Delta U=Q=m_1 c_1 \Delta t$ nhiệt độ lúc sau của hệ là $35,6^\circ C$.
		\end{itemchoice}
	}
\end{ex}

\begin{ex} \nguon{4}
	Để xác định nhiệt độ của một lò, người ta đưa vào lò một miếng sắt khối lượng $75 \ g$. Khi miếng sắt có nhiệt độ bằng nhiệt độ của lò, người ta lấy ra và thu ngay vào một nhiệt lượng kế có khối lượng $200 \ g$ có chứa $500 \ g$ nước ở nhiệt độ $30^\circ C$, khi hệ cân bằng nhiệt thì nhiệt độ của nước tăng lên đến $45^\circ C$. Coi nhiệt dung riêng của nước gần đúng là $4200 \ J/kg.K$, của sắt là $478 \ J/kg.K$ và của nhiệt lượng kế là $418 \ J/kg.K$. Giả sử hệ không trao đổi nhiệt với môi trường bên ngoài.
	\choiceTF[t]
	{Trong quá trình trên độ giảm nội năng của miếng sắt bằng độ tăng nội năng của nước}
	{\True Độ tăng nội năng của nước là $31500 \ J$}
	{\True Nhiệt độ của lò xấp xỉ là $958,6^\circ C$}
	{Miếng sắt đã truyền một nhiệt lượng là $35052 \ J$ cho nước và nhiệt lượng kế}
	\loigiai{
		\begin{itemchoice}
			\itemch Sai. Vì còn có nhiệt lượng truyền sang nhiệt lượng kế.
			\itemch Đúng
			\itemch Đúng
			\itemch Sai. Nhiệt lượng miếng sắt truyền cho nước và nhiệt lượng kế là $32754 \ J$.
		\end{itemchoice}	
	}
\end{ex}









